% !TeX program = xelatex
% Encoding: UTF-8
% 編譯方式：xelatex physics_math_W02_2026-10-09_dark.tex
\documentclass[12pt,a4paper]{article}
\usepackage[margin=2.3cm]{geometry}
\usepackage{fontspec}
\usepackage{xeCJK}
\setmainfont{Times New Roman}
% Windows 優先使用標楷體；其他平台依序選擇備援中文字型。
\IfFontExistsTF{標楷體}{
  \setCJKmainfont[AutoFakeBold=2]{標楷體}
}{
  \IfFontExistsTF{Noto Serif TC}{
    \setCJKmainfont{Noto Serif TC}
  }{
    \IfFontExistsTF{Microsoft JhengHei}{
      \setCJKmainfont{Microsoft JhengHei}
    }{
      \setCJKmainfont{Noto Sans TC}
    }
  }
}
\usepackage{amsmath,amssymb}
\usepackage{enumitem}
\usepackage{setspace}
\usepackage{xcolor}
\pagecolor{black}
\color{white}
\usepackage[unicode,hidelinks]{hyperref}
\setstretch{1.18}
\setlength{\parindent}{0pt}
\setlength{\parskip}{0.45em}
\setlist{itemsep=0.3em,topsep=0.3em}
\emergencystretch=2em
\title{物理數學短練習 W2\\對稱矩陣、位能與振動}
\author{}
\date{2026 年 10 月 9 日至 10 月 15 日}
\begin{document}
\maketitle
\vspace{-1.2em}
\textbf{對應計畫：}新版 Q1，W2：Strang 6-4 譜定理與正交對角化。
本次按計畫安排，不預設已完成前一週的診斷。

\textbf{時間與負荷：}共 3 題，約 35--40 分鐘，可分次完成。
輪班疲累時先做子題 1、2，子題 3 留作選做。
本份練習是一組大題，分成三個互有關聯的子題，連結線性代數、向量分析、微分方程與分析力學。
先完成題目，再翻至提示；完整解法另起頁。

\section{題目：一組大題、三個子題}
\subsection*{共用矩陣、符號與條件}
三個子題共用下列實對稱矩陣：
\[
    A=\begin{pmatrix}2&-1\\-1&2\end{pmatrix}.
\]
\begin{itemize}
    \item \(A\)：無因次矩陣，三個子題均使用同一個 \(A\)。
    \item \(Q\)、\(\Lambda\)：子題 1 求得的正交矩陣與對角矩陣；
          \(\mathbf e_1,\mathbf e_2\) 是對應的正交規範特徵向量。
    \item \(\mathbf r=(x,y)^T\)：位置向量，\(x,y\) 的單位為公尺。
    \item \(k>0\)：單位為 \(\mathrm{N/m}\)；\(m>0\)：質量，單位為公斤。
    \item \(a>0\)：初始位移尺度，單位為公尺；\(t\)：時間，單位為秒。
    \item 上點代表對時間微分；\(V\) 是位能，\(\mathbf F=-\nabla V\) 是力場。
\end{itemize}
\textbf{作答關聯：}子題 1 的特徵值與特徵向量可直接沿用至子題 3；
子題 2 使用共用矩陣 \(A\)，其力場可直接沿用至子題 3。
\subsection*{子題 1：正交對角化（約 12 分鐘，必做）}
使用上方的共用矩陣 \(A\)。
\begin{enumerate}
    \item 求全部特徵值與對應的特徵向量。
    \item 將特徵向量規範化，組成正交矩陣 \(Q\)，寫出
          \[
              A=Q\Lambda Q^T.
          \]
    \item 若 \(A\mathbf u=\lambda\mathbf u\)、\(A\mathbf v=\mu\mathbf v\)，
          且 \(\lambda\ne\mu\)，利用 \(A^T=A\) 證明 \(\mathbf u\) 與 \(\mathbf v\) 正交。
\end{enumerate}

\subsection*{子題 2：位能的梯度（約 8 分鐘，必做）}
\textbf{可沿用：}共用矩陣 \(A\)；如有需要，也可使用子題 1 的譜分解。

設位置向量為 \(\mathbf r=(x,y)^T\)，位能為
\[
    V(x,y)=\frac{k}{2}\mathbf r^TA\mathbf r,\qquad k>0.
\]
其中 \(x,y\) 的單位為公尺，\(k\) 的單位為 \(\mathrm{N/m}\)。
\begin{enumerate}
    \item 展開 \(V\)，求梯度 \(\nabla V\) 與力場 \(\mathbf F=-\nabla V\)。
    \item 求滿足 \(\nabla V=\mathbf0\) 的駐點，判斷它是否為嚴格極小值。
\end{enumerate}
所需定義：
\[
    \nabla V=
    \begin{pmatrix}
        \partial V/\partial x \\
        \partial V/\partial y
    \end{pmatrix}.
\]

\subsection*{子題 3：用特徵向量解振動方程（約 15 分鐘，選做）}
\textbf{可沿用：}子題 1 的特徵值、正交規範特徵向量，以及子題 2 的力場；不必重新計算。

質量為 \(m>0\) 的粒子在子題 2 的位能中運動：
\[
    m\ddot{\mathbf r}=-kA\mathbf r.
\]
初始條件為
\[
    \mathbf r(0)=\begin{pmatrix}a\\0\end{pmatrix},
    \qquad \dot{\mathbf r}(0)=\mathbf0,\qquad a>0.
\]
其中 \(m\) 的單位為公斤，\(a\) 的單位為公尺，時間 \(t\) 的單位為秒。
沿子題 1 的兩個正交規範特徵向量，寫成
\[
    \mathbf r(t)=c_1(t)\mathbf e_1+c_2(t)\mathbf e_2.
\]
求 \(c_1,c_2\)，以及 \(x(t),y(t)\)。

可使用：
\[
    \ddot c+\omega^2c=0
    \quad\Longrightarrow\quad
    c(t)=C\cos(\omega t)+D\sin(\omega t).
\]

\clearpage
\section{提示：需要時再看下一層}
\subsection*{子題 1}
\textbf{第一層：}先算 \(\det(A-\lambda I)\)。

\textbf{第二層：}試試 \((1,1)^T\) 與 \((1,-1)^T\)。
證明正交時，從兩種方式計算 \(\mathbf u^TA\mathbf v\)。

\subsection*{子題 2}
\textbf{第一層：}先把矩陣乘法展開，再取偏導數。

\textbf{第二層：}使用
\[
    x^2-xy+y^2=\frac12(x-y)^2+\frac12(x^2+y^2).
\]

\subsection*{子題 3}
\textbf{第一層：}對每個特徵向量使用 \(A\mathbf e_i=\lambda_i\mathbf e_i\)。

\textbf{第二層：}兩個係數各自滿足
\[
    \ddot c_i+\frac{k\lambda_i}{m}c_i=0.
\]
用 \(c_i(0)=\mathbf e_i^T\mathbf r(0)\) 求初始係數。

\clearpage
\section{完整解法}
\subsection*{子題 1}
特徵方程為
\[
    \det(A-\lambda I)
    =(2-\lambda)^2-1
    =(\lambda-1)(\lambda-3)=0.
\]
因此
\[
    \lambda_1=1,\qquad\lambda_2=3.
\]
代入 \(A\mathbf v=\lambda\mathbf v\)，得到
\[
    \lambda_1=1:\quad\mathbf v_1=\begin{pmatrix}1\\1\end{pmatrix},
    \qquad
    \lambda_2=3:\quad\mathbf v_2=\begin{pmatrix}1\\-1\end{pmatrix}.
\]
兩向量內積為零，長度皆為 \(\sqrt2\)，所以正交規範基為
\[
    \mathbf e_1=\frac1{\sqrt2}\begin{pmatrix}1\\1\end{pmatrix},
    \qquad
    \mathbf e_2=\frac1{\sqrt2}\begin{pmatrix}1\\-1\end{pmatrix}.
\]
取
\[
    Q=\frac1{\sqrt2}\begin{pmatrix}1&1\\1&-1\end{pmatrix},
    \qquad
    \Lambda=\begin{pmatrix}1&0\\0&3\end{pmatrix}.
\]
因為 \(Q^TQ=I\)，且 \(AQ=Q\Lambda\)，所以
\[
    \boxed{A=Q\Lambda Q^T}.
\]
也能寫成譜分解：
\[
    A=\mathbf e_1\mathbf e_1^T+3\mathbf e_2\mathbf e_2^T.
\]

\textbf{不同特徵值對應的特徵向量正交：}
由 \(A\mathbf v=\mu\mathbf v\)，
\[
    \mathbf u^TA\mathbf v=\mu\,\mathbf u^T\mathbf v.
\]
另一方面，由 \(A^T=A\)，
\[
    \mathbf u^TA\mathbf v
    =(A\mathbf u)^T\mathbf v
    =\lambda\,\mathbf u^T\mathbf v.
\]
相減得
\[
    (\lambda-\mu)\mathbf u^T\mathbf v=0.
\]
因 \(\lambda\ne\mu\)，故
\[
    \boxed{\mathbf u^T\mathbf v=0}.
\]
\textbf{注意：}特徵向量的正負號可以改，但 \(Q\) 中向量的順序必須和
\(\Lambda\) 中特徵值的順序一致。

\subsection*{子題 2}
先展開：
\[
    V=\frac{k}{2}(2x^2-2xy+2y^2)=k(x^2-xy+y^2).
\]
所以
\[
    \nabla V=k\begin{pmatrix}2x-y\\2y-x\end{pmatrix}=kA\mathbf r,
    \qquad
    \boxed{\mathbf F=\begin{pmatrix}-2kx+ky\\kx-2ky\end{pmatrix}}.
\]
求駐點：
\[
    2x-y=0,\qquad 2y-x=0
    \quad\Longrightarrow\quad
    \boxed{x=y=0}.
\]
為了判斷極值，改寫為
\[
    V=\frac{k}{2}\big[(x-y)^2+x^2+y^2\big].
\]
當 \((x,y)\ne(0,0)\) 時，上式嚴格大於零；原點的位能為零。
因此，\textbf{原點是嚴格的全域極小值}。

單位也吻合：
\[
    [V]=(\mathrm{N/m})\,\mathrm m^2=\mathrm J,\qquad
    [\nabla V]=\mathrm{J/m}=\mathrm N.
\]
物理意義：在梯度非零的地方，位能沿梯度方向增加得最快，
而力指向位能下降的方向。

\subsection*{子題 3}
代入 \(\mathbf r=c_1\mathbf e_1+c_2\mathbf e_2\) 與
\(A\mathbf e_i=\lambda_i\mathbf e_i\)，得到
\[
    m(\ddot c_1\mathbf e_1+\ddot c_2\mathbf e_2)
    =-k(c_1\mathbf e_1+3c_2\mathbf e_2).
\]
由兩向量線性獨立，
\[
    \ddot c_1+\frac{k}{m}c_1=0,\qquad
    \ddot c_2+\frac{3k}{m}c_2=0.
\]
因此兩個角頻率為
\[
    \omega_1=\sqrt{\frac{k}{m}},\qquad
    \omega_2=\sqrt{\frac{3k}{m}}.
\]
兩者單位皆為 \(\mathrm{s^{-1}}\)。

將初始位置投影到兩個方向：
\[
    c_1(0)=\mathbf e_1^T\mathbf r(0)=\frac a{\sqrt2},\qquad
    c_2(0)=\mathbf e_2^T\mathbf r(0)=\frac a{\sqrt2}.
\]
初始速度為零，故 \(\dot c_1(0)=\dot c_2(0)=0\)，所以
\[
    c_1(t)=\frac a{\sqrt2}\cos(\omega_1t),\qquad
    c_2(t)=\frac a{\sqrt2}\cos(\omega_2t).
\]
還原位置：
\[
    \boxed{x(t)=\frac a2[\cos(\omega_1t)+\cos(\omega_2t)]},
\]
\[
    \boxed{y(t)=\frac a2[\cos(\omega_1t)-\cos(\omega_2t)]}.
\]
檢查 \(t=0\)：得到 \(x(0)=a,y(0)=0\)；對時間微分後，
兩個初始速度也都是零。代回各係數的微分方程，即滿足原運動方程。

\textbf{物理意義：}特徵向量決定互不耦合的振動方向，
特徵值則決定各方向的恢復力強度。這正是正交對角化在力學中的用途。

\section{本週重點整理}
\begin{enumerate}
    \item \textbf{譜定理：}實對稱矩陣具有實特徵值，並能選出由特徵向量
          組成的正交規範基。重複特徵值的特徵空間內，仍需選取正交規範基。
    \item \textbf{別漏掉規範化：}只有互相垂直還不夠；要讓 \(Q^TQ=I\)，
          每個向量的長度也必須為 1。
    \item \textbf{銜接分析力學：}由
          \[
              L=\frac12m(\dot x^2+\dot y^2)-V
          \]
          使用 Euler--Lagrange 方程，可得到本題的
          \(m\ddot{\mathbf r}=-\nabla V\)。
    \item \textbf{銜接電磁學：}靜電場同樣由負梯度得到：
          \(\mathbf E=-\nabla\phi\)。注意 \(\phi\) 是電位；
          電荷 \(q_{\mathrm e}\) 的位能是 \(V=q_{\mathrm e}\phi\)。
    \item \textbf{銜接量子力學：}「求特徵值，再沿正交特徵向量分解」
          會延伸到量子力學的厄米算符。本週先掌握實對稱矩陣即可。
\end{enumerate}

\subsection*{自我檢核}
若子題 3 改成 \(\mathbf r(0)=a\mathbf e_1\)、初始速度仍為零，
運動中會出現幾個角頻率？請用「初始位置的投影」說明理由。
\end{document}